\section{Bombeo secundario seleccionado y separación hidráulica: ENV-37}
\label{sec:ambiente-env37}
\textbf{Solución elegida y frontera primaria.} Se seleccionan dos bombas simples Wilo \texttt{Stratos MAXO 65/0,5-16 PN6/10}, artículo \texttt{2164595}, una por cada circuito hidráulico independiente de frío. Se conserva la arquitectura de dos Carrier 30RB065-SS, uno de servicio y otro de reserva, con circuito, inercia, aislamiento y bombeo propios. Cada circuito secundario puede alimentar las cuatro rutas de tratamiento mediante cabeceras propias y selección hidráulica que impida unir simultáneamente los dos circuitos. El banco A y el B conservan sus dos rutas independientes de aire/serpentín; una ruta activa por banco y las otras disponibles. No se denomina bomba de reserva a dos bombas puestas en serie: la sustitución se realiza transfiriendo a la otra fuente/circuito completo, con continuidad de temperatura y caudal aún por verificar. Las válvulas de transferencia, antirretornos, aislamiento y elemento de desconexión por fuga permanecen por seleccionar; el esquema no acredita disponibilidad hidráulica final.

El secundario se desacopla del evaporador por un depósito de inercia de cuatro conexiones de cada circuito, conservando 1.500 L útiles proyectados. Su selección real y mezcla estratificada deberán demostrar entrega de agua 5 grados C en el peor relevo. La bomba seleccionada circula entre ese separador y los serpentines; no se le atribuyen las pérdidas del evaporador ni se reduce el caudal primario Carrier a la demanda instantánea de las válvulas. El dato nominal Carrier de 10,8 m$^3$/h no es caudal mínimo de evaporador verificado ni prueba de capacidad a agua 5/10 grados C. La bomba primaria y la curva corregida de fuente quedan abiertas; no se completa RT-06.13 por esta selección secundaria.

\textbf{Variante y curva publicada.} La ficha individual confirma DN65, distancia entre bridas 340 mm, alimentación monofásica 230 V $\pm10\,\%$/50--60 Hz, agua admisible $-10$ a $110\,^{\circ}$C, ambiente $-10$ a $40\,^{\circ}$C, IPX4D, potencia de entrada $P_1$ de 20 a 1.440 W y corriente de 0,3 a 6,23 A. Se selecciona la versión estándar 2164595, no la R7 2217962 ni PN16 2186283. La altura máxima individual publicada es 16,1 m; $Q_{max}=54,8$ m$^3$/h se refiere a altura próxima a cero y no es el caudal disponible a altura máxima. El catálogo primario de la misma variante presenta la envolvente de presión constante y variable: a 16,17 m$^3$/h la curva máxima se lee aproximadamente entre 15 y 16 m; se usa como cribado gráfico y no punto numérico garantizado. La altura de contraste de 12 m queda dentro de esa envolvente a 8,08 y 16,17 m$^3$/h; no se adopta como consigna local operativa. La solución elegida regula diferencial remoto en el ramal índice según se desarrolla más adelante. La selección numérica del fabricante debe confirmar caudal, altura, tolerancia, agua 5/10 y potencia real de cada estado antes de compra \parencite{lote37-wilo-producto}\parencite[p.~111]{lote37-wilo-curva}.

\textbf{Balance del secundario.} Se toma sólo como contraste el punto ideal de ENV-34 de frío 23,516928 kW por banco y agua 5/10 grados C: 4,041903 m$^3$/h por ruta, 8,083807 para dos rutas activas y 16,167614 para cuatro rutas en solape. La salida ideal saturada 7,3925 grados C utilizada en ese contraste no es una consigna adoptada ni rendimiento OEM; incertidumbre de instalación, bypass, ganancias de trazado y calor real de bomba pueden aumentar la demanda. La revisión se integra con el dimensionamiento húmedo antes de fijar caudal final.

Se conserva del contraste ENV-31 la rama de diámetro interior mínimo 40 mm, longitud ida/retorno 40 m, rugosidad 0,0015 mm, suma $K=12$, agua sin glicol con $\rho=999,8$ kg/m$^3$ y viscosidad 1,385 mPa s. Al nuevo caudal de contraste, su pérdida es 14,462599 kPa y la ZTVB25-8 completamente abierta pierde 25,526536 kPa. Con cota propia del resto de rama de 55 kPa, el ramal suma 80,526536 kPa, autoridad $25,526536/(25,526536+55)=0,316995$ y quedan 40,537401 kPa para serpentín y accesorios de rama no incluidos. Es condición de selección de esos equipos, no su curva recibida.

Para la distribución común secundaria se fija diámetro interior mínimo 80 mm, longitud total de ida/retorno máxima 50 m, rugosidad 0,0015 mm y suma $K=15$, incorporando en esos $K$ las transiciones inmediatas a DN65 de bomba. Son dimensiones de obra nueva. A 16,167614 m$^3$/h se obtiene $v=0,893457$ m/s, $Re=51.597$, factor Darcy por Haaland $f=0,0206132$ y pérdida común calculada 11,126899 kPa. No se aplican esos valores a una tubería comercial sin verificar diámetro efectivo, material y accesorios.
\[
\Delta p_{com}=\left(f\frac{50}{0,080}+15\right)\frac{999,8(0,893457)^2}{2}=11,126899\;\mathrm{kPa}.
\]
Con $g=9,80665$ m/s$^2$, la altura de diseño 12 m equivale a $999,8(9,80665)(12)/1000=117,656264$ kPa. Se limita a 20 kPa la pérdida adicional de separador, filtro común, dispositivos de transferencia/aislamiento y accesorios no contados; el total presupuestado es $80,526536+11,126899+20=111,653435$ kPa, dejando 6,002829 kPa respecto a la altura de contraste, que no es consigna operativa. Los 20 kPa son límite de selección, no catálogo de válvulas/filtro. Si las curvas reales, ensuciamiento o caudal final lo exceden, se recalculan distribución y conjunto; no se impone una altura que resulte insuficiente. La altura geométrica no se añade como pérdida continua de un circuito cerrado, pero sí determina llenado y presión estática.

\textbf{Cierre de ramas y presión.} La cota de altura máxima de catálogo 16,1 m equivale, para el agua proyectada, a 157,855488 kPa, menor que el diferencial máximo de 200 kPa de ZTVB25-8. Esto aporta un contraste aun si el modo se cambia a máxima velocidad: no se suman las alturas de dos bombas ni se introduce un primario en serie en el secundario. La regulación ordinaria elegida mantiene diferencial remoto en el ramal índice, hasta el límite de curva y la condición de parada sin flujo configurada. Con ramas cerradas se comprueba presión en todos los ramales; una medida remota no acredita por sí sola máximo en otra rama. La lectura no constituye prueba de transitorio, tolerancia de curva ni límite de seguridad mecánico. Se exige confirmar presión máxima de cierre/transferencia con curvas numéricas y medir sobrepresión; queda por seleccionar limitación independiente o corte si los picos superan la cota. Los 10 bar de PN son presión estática de cuerpo, distintos de 157,855 kPa de altura diferencial y de 200 kPa de válvula.

El manual fija presión mínima de aspiración manométrica de 0,7 bar para DN65 con altura máxima 16 m y agua entre $-10$ y $+50\,^{\circ}$C. Se establece una cota propia de presión en boca de aspiración de al menos 1,0 bar en todas las condiciones de obra, con vaso de expansión/presurización aún por seleccionar. Con 1,0 bar de aspiración y altura máxima nominal, la descarga sería aproximadamente 2,5786 bar, sin contar expansión o transitorio. No se toma esa cifra como presión máxima definitiva ni como NPSH recibido. Se verifican altura de llenado, presión en el punto alto, expansión, alivio y máxima presión de cada accesorio antes de ejecutar \parencite[p.~15]{lote37-wilo-instalacion}.

\textbf{Control y continuidad eléctrica.} Se elige regulación nativa de la bomba en modo ramal índice de diferencial constante remoto, con dos transmisores Wilo DDG20 4--20 mA, artículo 2136456, uno por circuito. El fabricante identifica conexión directa al AI1 o AI2 de Stratos MAXO y rango 0--2 bar, alimentación nominal 24 V, admisible 15--30 V DC, medio -20 a 80 grados C e IP55. El manual de conexión usa marrón3 positivo, azul2 referencia y negro1 señal; la referencia está unida a carcasa PELV y debe conciliarse con tierra, sin unir referencias A/B. Se fija demanda propia de 82,8 kPa en el ramal hidráulicamente más desfavorable: para escala 0--200 kPa corresponde señal 10,624 mA. Con cota del resto de rama55 kPa, la válvula dispone27,8 kPa, autoridad27,8/82,8=0,33575, y requiere Kv efectivo4,041903/sqrt(0,278)=7,6659, inferior aKvs8. No se convierte Kv en posición lineal de actuador sin curva instalada. Bajo solape, demanda de altura incluye82,8+11,126899+20=113,926899 kPa, altura11,61963 m aproximadamente; permanece dentro del cribado de curva, pendiente selección numérica. El ramal índice se verifica en cada estado; transferencia, cambio de ruta y ubicación de tomas deben demostrar que ninguna rama queda menos favorecida. No consume AO WAGO ni entradas AI del nodo ya ocupado. El DDG se conecta al AI de bomba; suministro activo24V del borne11 yGND12 sólo aparece tras configurar el AI y admite máximo50mA. La carga máxima admisible DDG es 500 ohm y el AI de bomba publica hasta 300 ohm para 4--20 mA; se corrige la falta de esa cota al cotejar ambos manuales. Consumo total DDG, exactitud, alimentación mínima y caída de cable aún requieren cierre: esa conexión permanece sin habilitar hasta verificarlos o seleccionar fuente compatible independiente. No se inventa exactitud ni potencia del DDG; la demanda eléctrica es adicional si no queda cubierta por P1 OEM de la bomba. Los 82,8 kPa son ajuste de diseño pendiente de incertidumbre de medida, no ajuste seguro ya validado \parencite{lote37-ddg-producto}\parencite[pp.~19--20, 23--24]{lote37-ddg-manual}\parencite[p.~36]{lote37-wilo-instalacion}. El SSM se configura como fallo colectivo y advertencias, con retardo de disparo propio previsto de cero segundos y lectura de retorno; el manual confirma que COM75/OK76 está cerrado tanto sin fallo como sin electricidad, por lo que ese contacto no detecta pérdida de alimentación. Se requiere supervisión de energía y el SBM como indicación de marcha conforme al manual; marcha tampoco equivale a caudal hidráulico de serpentín. Se requieren entradas/relés auxiliares independientes por seleccionar: las ocho DI por banco de ENV-34 ya están reservadas y no se duplican canales disponibles. El fallo retira disponibilidad de ese circuito y exige conmutación verificada del otro; no habilita recarga ni sustituye la cadena H$_2$ \parencite[pp.~18, 58--59]{lote37-wilo-usuario}.

Cada bomba se alimenta desde el camino AC protegido asociado a su fuente, con L/N/PE y seccionamiento propio, sin 400 V entre fases. Una activa reserva 1,440 kW y $230(6,23)/1000=1,43290$ kVA; en relevo de dos se reservan 2,880 kW y 2,86580 kVA. Son máximos de placa/ficha y sus productos para 230 V; la potencia aparente así calculada resulta ligeramente inferior a $P_1$ máximo publicado por redondeos/condiciones distintas, por lo que para el balance se fija cota propia de 1,60 kVA por bomba, 3,20 kVA de dos, sin presentar un factor de potencia calculado inválido. La recepción usa la corriente máxima en toda la banda de tensión y el par de valores simultáneos del fabricante para cerrar kW/kVA y pérdidas.

El manual declara picos de hasta 10 A al conectar alimentación, relé externo al menos 10 A/250 V AC, conexión fija/seccionamiento de todos los polos con separación mínima 3 mm, conductor admitido de red 3 por 1,5 a 3 por 2,5 mm$^2$ y fuga por bomba hasta 3,5 mA. La sección de obra y protección se calculan por longitud/caída/selectividad; ese intervalo de borne no es dimensionamiento del alimentador. Por $P_1>200$ W puede aparecer corriente residual DC o pulsante; la selección de diferencial se verifica con normativa aplicable y fabricante, sin inferir tipo A suficiente. Se evita control de red por triac y se respeta el límite de cien conexiones diarias por red publicado. Arranque, UPS, autonomía y reparto de fases quedan por conciliar con esta carga adicional \parencite[pp.~30--33]{lote37-wilo-instalacion}.

Como cota térmica conservadora se agrega a la selección de fuente hasta 1,440 kW por bomba activa, sin garantizar que todo $P_1$ termine en agua. Dos bancos al punto ideal requieren al menos $2(23,516928)+1,440=48,473856$ kW antes de ganancias de trazado, incertidumbre/serpentín, bomba primaria y otros accesorios. El solape de cuatro rutas a plena demanda del contraste exige 94,067712 kW de frío de proceso, antes de calor de bombeo, y supera los 63 kW nominales de un Carrier existente en la selección de arquitectura; se prohíbe atribuir capacidad N+1 a cuatro rutas húmedas simultáneas a plena carga por la sola curva de caudal. Se conserva solape de aire, pero el agua de reserva se habilita por etapas con reparto verificable de carga o se sustituye fuente; esa secuencia térmica todavía no está validada. La bomba capaz de 16,167614 m$^3$/h no demuestra refrigeración para esa demanda.

\textbf{Montaje y condensación.} La bomba individual pesa aproximadamente 31,6 kg netos; su envolvente declarada es 485 por 340 por 289 mm y requiere eje horizontal, dirección de flujo de carcasa, soporte independiente y conexiones sin esfuerzo. Se reservan dentro de cada módulo hidráulico vigente de 2 por 2 m una bancada de obra de 1,2 por 0,8 m y acceso frontal propio de 0,8 m. La reserva no acredita encaje de depósito/primario/válvulas; se fija sólo para bomba secundaria y accesorios inmediatos. No se modifica el predio ni rutas de escape \parencite{lote37-wilo-producto}\parencite[pp.~23--24]{lote37-wilo-instalacion}.

Se seleccionan dos carcasas aislantes de agua fría Wilo \texttt{ClimaForm Stratos MAXO 65/0,5-12/16}, artículo \texttt{2201740}, y dos kits de contrabridas PN10 DN65 \texttt{2105583}; el catálogo incluye dos contrabridas de acero soldables EN1092-1 tipo 11 con juntas, tornillos y tuercas por kit. No se usa en agua fría la carcasa térmica estándar incluida para calefacción: el manual la limita a agua superior a 20 grados C. ClimaForm se individualiza como accesorio compatible, sin atribuir conductividad, espesor o clase al fuego no publicados. La recepción debe demostrar su superficie por encima del rocío local; no sustituye los 100 mm de reserva exterior INT-34 en tubos/puentes. Sólo se aísla el cuerpo hasta separación del motor, dejando libres canales de condensado; motor/electrónica no se cubren. La carcasa/accesorio queda dentro de protección no condensante aún por dimensionar, fuera de gas de banco. El manual requiere impedir condensación electrónica y contempla marcha continua o calefacción de acompañamiento: la bomba de reserva detenida requiere resolver esa medida con potencia y control propios antes de afirmar disponibilidad, sin hacerla girar contra circuito cerrado \parencite[p.~111]{lote37-wilo-curva}\parencite[p.~29]{lote37-wilo-instalacion}.

Los cuatro drenajes de aire limpio del desarrollo integrado conservan caudal de condensado de contraste de al menos 22,203622 kg/h por ruta, hasta 88,814488 kg/h para cuatro, sin mezclarlo con retorno H$_2$ ni con evacuación del motor de bomba. El serpentín real puede exigir más por condición/factor de bypass; se recalcula el drenaje al seleccionarlo. Esta selección concreta bomba secundaria, curva de cribado, distribución y cantidades. RT-06.03, RT-06.09 y RT-06.13 siguen pendientes por montaje/circuito primario, fuente y serpentín corregidos, capacidad N+1/solape térmico, protección/continuidad, accesorios hidráulicos completos, condensación y seguridad independiente.

\textbf{Frontera de pérdida de camino.} ENV-40 reemplaza la alimentación única del riel ambiental de cada banco por desacoplamiento A/B seleccionado. El mando permanece en una CPU por banco; la continuidad de los receptores AC y del tratamiento completo conserva la brecha siguiente. Transferir agua al circuito de la otra fuente no garantiza mando de las válvulas, impulsión y vigilancia del banco cuyo camino se retiró. La regulación remota de bomba no sustituye ese mando; seleccionar dos bombas no acredita tratamiento N+1 completo. El esquema final debe demostrar alimentación/control de las rutas que sostienen ambos bancos en cada retirada y su estado seguro de falla, con dispositivos seleccionados, potencia y maniobras incluidos. Hasta resolver esa dependencia se conserva pendiente continuidad de climatización y no se declara disponible una ruta por tener sólo agua presurizada.

\textbf{Comprobación de carga de DDG por el integrador.} El manual DDG fija alimentación 24 V DC ±20 por ciento (19,2--28,8 V) y carga máxima 500 $\Omega$; el manual de instalación de bomba p.~36 fija carga AI 4--20 mA no mayor de 300 $\Omega$. Se proyecta por sensor máximo 10 m de ida y conductor efectivo mínimo 0,5 mm$^2$, con resistividad propia 0,0175 $\Omega$ mm$^2$/m: ida/retorno 20 m añaden 0,7 $\Omega$ antes de temperatura/bornes. Reservando 1 $\Omega$ completo de conexión y cable, $R_{total}\le301~\Omega<500~\Omega$. Esto comprueba compatibilidad de carga bajo ese presupuesto, sin dar una ficha comercial al cable. Deben verificarse sección, temperatura y tendido reales, tensión de sensor en transferencia y consumo completo dentro de los 50 mA de la salida activa. Sobrecarga/corto de 24 V hace fallar todas las funciones de entrada de bomba según su manual; esa causa retira disponibilidad y conserva investigación, sin considerar la presión de referencia congelada como medición válida. Exactitud/latencia, alimentación y respuesta de falla todavía impiden habilitar regulación integral.\parencites[p.~19]{lote37-ddg-manual}[p.~36]{lote37-wilo-instalacion}.
